% Note: part labels below are those of the official solution ((1.1)-(3.5)).
% The statement transcribed in the book is the released exam version, whose
% numbering (1.1.1, 1.2.1, ...) differs slightly and which omits the
% distance part (2.4) and the whole section 3 (speed of the gravitational
% wave); the full official solution is transcribed here.

\noindent\textbf{1\quad Data from the graph}
\begin{description}
\item[(1.1)] To make the scaling as accurate as possible, we take the time
  values far apart. When $x = 0.0$\,cm $t = -30$\,s, using the ruler we can
  also measure that at $t = 0$\,s we have $x = 11.1$\,cm. \hfill(1 point)

  This leads to
  \begin{align*}
    0.0\,\mathrm{cm} &= (-30\,\mathrm{s})A + B \\
    11.1\,\mathrm{cm} &= (0\,\mathrm{s})A + B = B \\
    A &= \frac{11.1\,\mathrm{cm}}{30\,\mathrm{s}} = 0.37\,\mathrm{cm/s} \\
    \therefore t &= \frac{x - 11.1}{0.37}\,\mathrm{s}
  \end{align*}
  \hfill(1 point)

\item[(1.2)] Analogues to the previous scaling,
  \begin{align*}
    0.0\,\mathrm{cm} &= C\log(30\,\mathrm{Hz}) + D \\
    5.3\,\mathrm{cm} &= C\log(500\,\mathrm{Hz}) + D
  \end{align*}
  \hfill(1 point)
  \begin{align*}
    5.3\,\mathrm{cm} &= C\log\left(\frac{500}{30}\right) \\
    C &= 4.34\,\mathrm{cm} \\
    D &= -6.41\,\mathrm{cm} \\
    \therefore f &= 10^{\left(\frac{y+6.41}{4.34}\right)}\,\mathrm{Hz}
  \end{align*}
  \hfill(2 points)

\item[(1.3)] It is not necessary to use $X$-scaling as students can choose
  convenient points from the graph.

  \begin{center}
  \begin{tabular}{ccccc}
    \hline
    $x$ (cm) & $y$ (cm) & $t$ (s) & $f$ (Hz) & $f^{-8/3}$ ($\mathrm{Hz}^{-8/3}$) \\
    \hline
    - & 0.5 & $-28.0$ & 39.1 & 5.67E-05 \\
    - & 0.5 & $-26.0$ & 39.1 & 5.67E-05 \\
    - & 0.6 & $-24.0$ & 41.3 & 4.92E-05 \\
    - & 0.6 & $-22.0$ & 41.3 & 4.92E-05 \\
    - & 0.7 & $-20.0$ & 43.5 & 4.27E-05 \\
    - & 0.8 & $-18.0$ & 45.9 & 3.71E-05 \\
    - & 0.9 & $-16.0$ & 48.4 & 3.22E-05 \\
    - & 1.0 & $-14.0$ & 51.0 & 2.79E-05 \\
    - & 1.1 & $-12.0$ & 53.8 & 2.43E-05 \\
    - & 1.2 & $-10.0$ & 56.7 & 2.11E-05 \\
    - & 1.4 & $-8.0$  & 63.1 & 1.59E-05 \\
    - & 1.6 & $-6.0$  & 70.1 & 1.20E-05 \\
    - & 1.8 & $-4.0$  & 78.0 & 0.90E-05 \\
    - & 2.3 & $-2.0$  & 101.7 & 0.43E-05 \\
    - & 3.4 & 0.0     & 182.4 & 0.09E-05 \\
    \hline
  \end{tabular}
  \end{center}
\end{description}

\noindent\textbf{2\quad Calculate system parameters}
\begin{description}
\item[(2.1)]
  \begin{align*}
    \frac{96}{5}\pi^{8/3}
      \left(\frac{G M_{chirp}}{c^3}\right)^{5/3} &= f^{-11/3}\dot{f} \\
    \therefore \left[\frac{96}{5}\pi^{8/3}
      \left(\frac{G M_{chirp}}{c^3}\right)^{5/3}\right]t
      + \left[\frac{3}{8}\right] f^{-8/3} + C &= 0
  \end{align*}
  \hfill(3 points)

\item[(2.2)]
  % FIGURE NEEDED: linear fit of f^{-8/3} (0 to 6e-5 Hz^{-8/3}) versus time
  % (-30 s to 0 s), 15 points on a descending straight line with fitted line
  % f(x) = -2.10E-06 x - 8.35E-08, R^2 = 9.94E-01; official 2022 DA
  % solutions, p. 3.
  Calculation of $f^{-8/3}$ values \hfill(3 points)

  Properly drawn graph \hfill(6 points)
  \begin{align*}
    \mathrm{slope} &= \frac{8}{3}\left[\frac{96}{5}\pi^{8/3}
      \left(\frac{G M_{chirp}}{c^3}\right)^{5/3}\right] \\
    M_{chirp} &= \left[\left(\frac{5}{256}\right)^{3/5}
      \frac{c^3}{G}\,\pi^{-8/5}\right](\mathrm{slope})^{3/5}
  \end{align*}
  \hfill(2 points)
  \[ \mathrm{slope} = 2.10 \times 10^{-6} \]
  \hfill(2 points)
  \[ \therefore M_{chirp} = 2.39 \times 10^{30}\,\mathrm{kg} = 1.20 M_\odot \]
  \hfill(2 points)

  with 95\% confidence interval
  \begin{align*}
    \Delta\mathrm{slope} &= 9 \times 10^{-8} \\
    \frac{\Delta M_{chirp}}{M_{chirp}}
      &= \frac{3}{5}\,\frac{\Delta\mathrm{slope}}{\mathrm{slope}} = 0.04 \\
    \Delta M_{chirp} &= 0.03 M_\odot
  \end{align*}

\item[(2.3)]
  \[ M_{chirp} = \frac{M_{chirp\,\mathrm{detector}}}{1+z}
     = \frac{1.20 M_\odot}{1.009783} = 1.19 M_\odot \]
  \hfill(2 points)

\item[(2.4)] According to the Hubble law,
  \[ D = \frac{zc}{H} = \frac{0.009783 \times 3 \times 10^5}{70}
     = 41.9\,\mathrm{Mpc} \]
  \hfill(2 points)

\item[(2.5)]
  \begin{align*}
    M_{chirp} &= \frac{(m_1 m_2)^{3/5}}{(m_1+m_2)^{1/5}} \\
    \therefore M_{chirp} &= m_2 \sqrt[5]{\frac{q^3}{1+q}}
  \end{align*}
  \hfill(2 points)

  This is an increasing function of $q$
  \begin{align*}
    \therefore m_{2,min} &= M_{chirp} 2^{1/5} = 1.37 M_\odot \\
    m_{2,max} &= 1.60 M_\odot
  \end{align*}
  \hfill(1 point)

  Also,
  \[ M_{chirp} = \frac{m_1}{q}\sqrt[5]{\frac{q^3}{1+q}}
     = \frac{m_1}{\sqrt[5]{q^2(1+q)}} \]
  \hfill(2 points)
  \begin{align*}
    \therefore m_{1,max} &= M_{chirp} 2^{1/5} = 1.37 M_\odot \\
    m_{2,min} &= 1.17 M_\odot
    % sic: the official solution labels this second value m_{2,min};
    % it is the minimum of the primary mass, m_{1,min}.
  \end{align*}
  \hfill(1 point)
\end{description}

\noindent\textbf{3\quad Speed of the gravitational wave}
\begin{description}
\item[(3.1)]
  % FIGURE NEEDED: plot of Fermi GBM event count (-100 to 200 counts/s)
  % versus time (0 to 4 s) from the table in the problem, points joined by
  % line segments, peak of ~179 at t ~ 2.4 s; official 2022 DA solutions,
  % p. 5.

\item[(3.2)] It appears that the event is starting from $\Delta t = 1.8$\,s.
  Some students may also write $\Delta t = 2.2$\,s as that is when the signal
  is truly above fluctuations seen before the event. We will accept both the
  answers. \hfill(1 point)

\item[(3.3)] By simple substitution we can calculate that
  \[ \frac{\Delta v}{v_{EM}} = \frac{\Delta t\, v_{GW}}{D}
     \approx \frac{\Delta t\, c}{D} \]
  \hfill(1 point)

\item[(3.4)]
  \[ \frac{\Delta v}{v_{EM}} = 4.2 \times 10^{-16}
     \text{ to } 5.1 \times 10^{-16} \]
  \hfill(1 point)

\item[(3.5)]
  \[ \Delta t' = \Delta t - 10 = -8.2\,\mathrm{s} \text{ to } 7.8\,\mathrm{s} \]
  % sic: the official solution omits the minus sign; the upper value should
  % be -7.8 s (= 2.2 s - 10 s).
  From this time delay
  \[ \frac{\Delta v}{v_{EM}} = -1.9 \times 10^{-15}
     \text{ to } 1.8 \times 10^{-15} \]
  % sic: correspondingly, the upper value should be -1.8e-15.
  \hfill(1 point)
\end{description}
